%% =====================================================================
%%  B7-T-06-02 -- what B7 contributes to "The Two-Tellings Test"
%%  -------------------------------------------------------------------
%%  LAYER A: APPEND-ONLY. Written once at the entry's write, then left
%%  alone. If the source has more to say later, add a bullet -- do not
%%  rewrite. Material found ONLY here goes in the `unique` slot with
%%  the unique flag set, which makes it undeletable (rule R10).
%% =====================================================================
%% @kind: source
%% @id: B7-T-06-02
%% @book: B7
%% @topic: T-06-02
%% @locator: ch. 5 (pp. 37--46), additive
%% @status: distilled
%% @unique: no
%% @integrated: yes
%% @updated: 2026-09-19

\dossier{B7}{T-06-02}{\bookshort{B7}}{ch. 5 (pp. 37--46), additive}

\begin{distilled}
  \item B7's verbal catalogue crosses the constructed-account family at several points: inconsistent statements (Montaigne's epigraph --- a bad memory is the enemy of lying); failure to answer the question asked; reluctance or refusal to answer; and the referral statement that answers every fresh question with some earlier occasion.
  \item Selective memory ('not that I recall', 'to the best of my knowledge') is the constructed account's preferred hedge: it makes every future inconsistency unfalsifiable at the time it is spoken.
  \item The constructed account, B7 notes, tends to be either evasively vague or overly specific --- and the two extremes can appear in the same speaker.
\end{distilled}
